\section{Discussion}
\label{sec:discussion}

\paragraph{What the factorization shows.}
A per-input router can fail because it cannot tell \emph{whether} the views disagree, \emph{which} wins when they do, how the latter shifts across domains, or because the apparent headroom is not real; measured separately, these have different weights here. Disagreement is moderately predictable (M1). Transfer matters in specific cells: for XLM-R on Algerian and Moroccan text, source $\hat\rho$ ranks direction at or below chance, and fitting it on the target turns a loss of up to 3.3 points into a small gain. But even fitted on the target, $\hat\rho$ gains at most 2.8 points over the best view in any cell (M2, M3). In our measurements the main limit is that the decision-point features say little about which view is right, even in domain, and a sizeable part of the oracle headroom is of the size that two classifiers differing only in their training sample also show (M4; a diagnostic reference, not a decomposition). The richer variants we tried did not change this (16 embedding components: $-1.2$ to $+1.1$ points, no consistent sign; gradient boosting, mE5 only: $-0.6$ to $+0.7$). These are limits of the features and estimators we evaluated; features computed after the second pass, fine-tuned representations or other estimators might do better, and our results do not show that the conditional advantage is intrinsically unpredictable. The same measurements can be made for other two-view choices with a frozen model, such as test-time augmentation or translate-test; we have not done so. Where oracle headroom is reported, a bootstrap-classifier reference helps to judge how much of it is classifier variability.

\paragraph{Where the signal is.}
The win rate on disagreements varies far more between encoder--dialect pairs (33--77\%) than our features resolve within a set, and on MYC every linguistic condition changes sign with the encoder. The useful decision is thus mostly set-level, which few labels resolve: with 64--128 target labels a set-level choice is on average within 0.3--0.7 points of SelectNorm with the same labels (descriptive). Targeted labels estimate the shift more efficiently (\cref{prop:efficiency}) but cannot fix a weak ranking.

\paragraph{Practical recommendations.}
For Arabizi with a frozen encoder, a linear head trained on Arabic script, and binary sentiment, our results suggest: (i) invest in normalizer safeguards first, since they matter more than routing (\cref{tab:ablation}); (ii) choose between normalizing and not at the level of a corpus or dialect, with on the order of a hundred labels; (iii) do not expect per-input routing that reads the first pass to save compute, except relative to running both views; (iv) when the classifier is trained on both scripts, learned routers, confidence thresholding or an ensemble sometimes gain a little over the better view, but can also lose on MYC.

\paragraph{Relation to prior findings.}
That transliteration helps unseen scripts on average is known \citep{muller2021unseen,moosa2023transliteration}; our results agree for Tunisian and Algerian Arabizi and add that the sign flips with the encoder on Moroccan text. Loss predictors for test-time transformations \citep{kim2020learningloss} and post-hoc deferral \citep{jitkrittum2023cascade} are the closest methods; SelectNorm is within one point of the direct-gain router in every cell, and the factorization adds a diagnosis, not an accuracy gain.
